% !TEX root = ../main.tex

\begin{abstract}[zh,title={摘\quad 要}]
  具身智能的发展正在推动机器人从实验室走向复杂开放的真实场景。视觉目标导航作为连接感知与决策的核心环节，要求机器人在弱先验条件下根据图像目标完成自主移动。近年来，目标掩码扩散导航模型（Goal Masked Diffusion Policies for Navigation and Exploration, NoMaD）等通用导航模型将扩散策略引入视觉导航，通过条件扩散过程生成多模态局部轨迹，展现了较强的表达能力。但现有验证大多围绕轮式平台展开，面向四足机器人的部署研究仍相对有限。

  本文以 NoMaD 为高层视觉导航核心，以云深处 Lite3 四足机器人为目标平台，围绕``基于动作扩散策略的四足机器人视觉目标导航''展开研究。主要工作包含三方面：第一，在离线层面对 NoMaD 进行部署导向优化，研究了去噪扩散隐式模型（Denoising Diffusion Implicit Model, DDIM）推理加速、无分类器引导（Classifier-Free Guidance, CFG）增强、测试时扩展（Test-Time Scaling, TTS）推理时搜索与视觉编码器升级，建立了统一的离线统计实验链路；第二，构建了由统一推理模块、状态机、导航主机、中层控制映射与平台接口组成的分层闭环系统，并在多关节接触动力学仿真器（Multi-Joint dynamics with Contact, MuJoCo）中完成了 Lite3 四足导航、探索与多目标任务的系统级验证；第三，形成了面向 Jetson Orin 与 Lite3 的真机部署方案，包括分阶段联调流程与多级安全机制。实验结果表明，DDIM-2 配置在保持轨迹质量的同时实现约 4.84 倍加速，MuJoCo 闭环验证中所有导航任务均成功到达目标，Orin 阶段一独立调试已通过验证。本文将原本面向轮式平台的 NoMaD 扩展为面向四足平台的可分析、可仿真、可部署的完整研究链路。
\end{abstract}

\begin{abstract}[en,title=ABSTRACT]
  The development of embodied intelligence is driving robots from laboratory settings into complex open-world scenarios. Visual goal navigation, as a core component connecting perception and decision-making, requires robots to achieve autonomous movement toward image-specified goals under weak prior conditions. Recently, general navigation models such as Goal Masked Diffusion Policies for Navigation and Exploration (NoMaD) have introduced diffusion policies into visual navigation, generating multi-modal local trajectories through conditional diffusion processes with strong expressiveness. However, existing validations primarily focus on wheeled platforms, and deployment research for quadruped robots remains limited.

  This thesis takes NoMaD as the high-level visual navigation core and the DeepRobotics Lite3 quadruped robot as the target platform, conducting research on ``visual goal navigation for quadruped robots based on action diffusion policy.'' The main contributions span three aspects. First, deployment-oriented optimizations are performed on NoMaD at the offline level, including Denoising Diffusion Implicit Model (DDIM) inference acceleration, Classifier-Free Guidance (CFG) enhancement, Test-Time Scaling (TTS) search, and visual encoder upgrades, establishing a unified offline statistical experiment pipeline. Second, a hierarchical closed-loop system comprising a unified inference module, state machine, navigation host, mid-level control mapping, and platform interface is constructed, with system-level verification completed in Multi-Joint dynamics with Contact (MuJoCo) simulation for Lite3 quadruped navigation, exploration, and multi-goal tasks. Third, a real-robot deployment scheme targeting Jetson Orin and Lite3 is developed, including staged commissioning procedures and multi-level safety mechanisms. Experimental results show that the DDIM-2 configuration achieves approximately 4.84$\times$ speedup while maintaining trajectory quality, all navigation tasks in MuJoCo closed-loop verification successfully reach their goals, and Orin Stage-1 standalone validation has passed. This work extends NoMaD from a wheeled-platform framework to an analyzable, simulatable, and deployable research pipeline for quadruped platforms.
\end{abstract}
